\chapter{The Middle Class: Structure, Ethos and the New Woman}

\section{Cultural vs Structural Change in Caste}

\begin{KeyConcept}
\begin{itemize}
\item For 'changes in caste since independence', distinguish \textbf{structural} from \textbf{cultural} change:
\item \textbf{Structural change} --- caste assumed a \textbf{new (political) avatar} (a new limb was added): \textbf{structural differentiation} --- from caste associations to multiple associations (CII, the National Women's Commission --- not caste-based). \textbf{Politicization of caste = structural change}.
\item \textbf{Cultural change} --- the old limbs are weakening: \textbf{secularization of caste = cultural change}. \textbf{Endogamy/hypergamy} is weakening (from restricted hypergamy to open intercaste marriage) and the \textbf{caste panchayat's} regulatory role is declining.
\end{itemize}
\end{KeyConcept}

\section{Obituary to Caste and the Peculiar Tenacity of Caste}

\begin{KeyConcept}
\begin{itemize}
\item \textbf{M. N. Srinivas's 'Obituary to Caste'} (an article, later the book \textbf{'Caste: Its Twentieth Century Avatar'}, c. 1990): he invited scholars to write on the changing role of caste. His follower \textbf{Jayaram} wrote of the \textbf{decline of caste as a system and its transformation into interest groups} --- decline due to \textbf{economic change} (occupation no longer linked to caste), transformation into interest groups due to \textbf{political change}.
\item Srinivas's line: \textbf{'caste as a system which has endured for over 2,000 years is dead or dying, although individual castes are thriving'} --- caste is now only \textbf{identity}, not a system. ('Caste is only identity now' --- the line by \textbf{Dipankar Gupta}, inspired by Srinivas.)
\end{itemize}
\end{KeyConcept}

\begin{KeyConcept}
\begin{itemize}
\item \textbf{Srinivas \& Beteille's article 'The Peculiar Tenacity of Caste'}: caste is in \textbf{decline in three areas} (religion/ritual, marriage, occupation) but \textbf{rising in one (politics)}.
\item \textbf{Religion/ritual}: earlier there was physical distancing by caste (in South India, the \textbf{Nairs} could approach within 16 feet of the \textbf{Namboodiri} Brahmins, other castes 64 or 128 feet); now there is a \textbf{secular trend of decline in the ritual basis of caste} --- but untouchability continues in a \textbf{new form: atrocity} (Rohith Vemula, mob/cow lynching).
\item \textbf{Marriage}: earlier endogamy + restricted hypergamy; now \textbf{intercaste marriage irrespective of anuloma/pratiloma}. Beteille: the \textbf{Radhi} and \textbf{Barendra} (both Bengali Brahmin jatis) earlier could not intermarry; today that is merely a marriage within the Brahmin community. Intercaste marriage happens \textbf{within adjacent sub-castes} (Brahmin--Baidya, Baidya--Kayastha). There is \textbf{'role distance'} (Goffman) --- 'I don't believe in caste, but by the way, what is your caste?'
\item \textbf{Occupation}: at the \textbf{general level} the caste--occupation link is weakening (mobility, merit, reservation); at the \textbf{specific level} it persists --- the \textbf{Teli} (oil-presser) sub-castes still use different traditional techniques (one bullock vs two bullocks).
\end{itemize}
\end{KeyConcept}

\section{European vs Indian Middle Class (Ahmad \& Rehfeld)}

\begin{KeyConcept}
\begin{itemize}
\item A comparison between the \textbf{European middle class (EMC)} and the \textbf{Indian middle class (IMC)}:
\item \textbf{Origin}: the EMC is \textbf{orthogenetic} (born within Europe, out of European conditions); the IMC is \textbf{heterogenetic} (born of colonialism).
\item \textbf{Function}: the EMC developed \textbf{to question authority} (the 'third estate' --- intellectuals, lawyers, doctors --- overthrew the first and second estates, the French Revolution); the IMC \textbf{initially supported British rule, then gradually questioned it} (moderates $\rightarrow$ extremists $\rightarrow$ revolutionaries $\rightarrow$ Gandhi), and became the \textbf{leader of the nationalist movement}.
\end{itemize}
\end{KeyConcept}

\section{Yogendra Singh: Urban vs Rural Middle Class}

\begin{KeyConcept}
\begin{itemize}
\item \textbf{Yogendra Singh}: the Indian middle class is \textbf{not homogeneous} --- it is divided into the \textbf{urban middle class (UMC)} and the \textbf{rural middle class (RMC)}.
\item \textbf{Similarity}: both share the same \textbf{ideological ethos} --- \textbf{conservative and utilitarian}.
\item \textbf{Difference --- antagonism}: the RMC shares an \textbf{antagonistic relation} with the UMC (entrepreneurs and professional groups). Reasons: \textbf{stagnation in agricultural development, unfavourable price policy, land fragmentation, and non-availability of non-farm occupations} --- so even the rich RMC faced \textbf{downward mobility} and felt alienated from the UMC (e.g. the farmer protests --- the RMC resenting the urban middle class).
\end{itemize}
\end{KeyConcept}

\section{Dipankar Gupta and Beteille: Mistaken Modernity and Middle Classes}

\begin{CriticalPoint}
\textbf{Dipankar Gupta} (in \textbf{'Mistaken Modernity'}): the Indian middle class is \textbf{only externally modern} --- it equates modernity with the \textbf{consumption of durable goods}. '\textbf{The bigger car they have, the more laws they are going to violate}' (Ferrari/Lamborghini owners who do not follow traffic rules). \textbf{Sitaram Yechury} calls this \textbf{'modern noxity'} (toxic modernity), not modernity.
\end{CriticalPoint}

\begin{KeyConcept}
\begin{itemize}
\item \textbf{Beteille}: India has \textbf{middle classes, not a middle class} --- the \textbf{upper middle class} (which follows a \textbf{consumerist culture} --- replacing goods with each new launch) and the \textbf{lower middle class} (only consumers, buying \textbf{cheaper imitations} of designer labels, e.g. at Palika Bazaar). The UMC looks up to the \textbf{Western elites}; the LMC looks up to the \textbf{UMC}.
\item \textbf{Material progress is taking place, but consumerism is yet to evolve} --- and the middle class must also have a \textbf{modern ethos} (which it lacks). Therefore the middle class is \textbf{neither traditional nor modern --- it is transitional} (consumption without ethos; a confused ethos; dependent on the market, not the state).
\item The middle class is \textbf{highly differentiated} and cannot be called homogeneous.
\end{itemize}
\end{KeyConcept}

\section{Women in the Middle Class and the New Woman}

\begin{KeyConcept}
\begin{itemize}
\item \textbf{Women in the middle class}: in the \textbf{upper middle class}, the daughter is educated \textbf{at par with the son} but \textbf{discouraged from higher education} (and passed off in marriage --- education can become a hindrance to finding a match). In the \textbf{lower middle class}, education is seen as a \textbf{passport to marriage} (the husband needs an educated wife). Tolerance for \textbf{inter-caste/inter-religious marriage is high in the UMC and low in the LMC} (norms override individuals).
\item \textbf{The new (middle-class) woman} --- empowered and ambitious: she will not sit silent, reports harassment, tolerates no gender inequality, and takes \textbf{instrumental roles} (CRPF Cobra commandant, Air Force pilot, army); she prefers IPS to IAS. She is the living symbol of the ongoing feminist wave.
\end{itemize}
\end{KeyConcept}

\begin{KeyConcept}
\begin{itemize}
\item \textbf{Conclusion}: the middle class has witnessed tremendous changes --- from the \textbf{Nehruvian civil-service-oriented salariate} to the \textbf{entrepreneur, IT and private-sector oriented} middle class; 'the middle class is no more just a class, it is a \textbf{set of classes}'.
\item Other middle classes: the \textbf{Dalit middle class} (Dalit consciousness and assertion --- studied by \textbf{Gurram Srinivas} of JNU) and the \textbf{Muslim middle class} (which supports reformist legislation like the abolition of triple talaq).
\end{itemize}
\end{KeyConcept}

\begin{QuickRevision}
\begin{itemize}
\item Structural change (politicization, new political avatar) vs cultural change (secularization, endogamy/caste panchayat weakening).
\item Srinivas's 'Obituary to Caste' (caste as a system is dead/dying; individual castes thrive; caste = only identity); Jayaram (decline $\rightarrow$ interest groups).
\item 'Peculiar Tenacity of Caste' (Srinivas \& Beteille): decline in religion/ritual (physical distancing $\rightarrow$ atrocity), marriage (restricted hypergamy $\rightarrow$ open intercaste within adjacent sub-castes; Radhi/Barendra; role distance), occupation (general decline, specific persistence --- Teli); rise only in politics.
\item Ahmad \& Rehfeld: EMC (orthogenetic, questions authority) vs IMC (heterogenetic, colonialism, led nationalism).
\item Yogendra Singh: urban vs rural middle class (same ethos; antagonism --- price policy, land fragmentation, no non-farm jobs, downward mobility).
\item Gurcharan Das: NEP 1991 $\rightarrow$ entrepreneurial surge; Dipankar Gupta ('Mistaken Modernity' --- externally modern, 'bigger car, more law-breaking'; Yechury's 'modern noxity'); Beteille (middle classes; transitional; consumerism without ethos).
\item Women: UMC (par education, no higher education) vs LMC (education = passport to marriage); the new woman (ambitious, instrumental roles).
\item Dalit middle class (Gurram Srinivas), Muslim middle class (reformist); 'a set of classes'.
\end{itemize}
\end{QuickRevision}

